% Part: computability
% Chapter: computability-theory
% Section: def-functions-self-reference

\documentclass[../../../include/open-logic-section]{subfiles}

\begin{document}

\olfileid{cmp}{thy}{slf}
\olsection{Netepake Fungsi Nganggo Rujukan Marang Awake Dhewe}

Lumrahe migunani yen kita bisa netepake fungsi
adhedhasar fungsi iku dhewe. Tuladhane, yen diwenehi
fungsi komputabel $k$, $l$, lan~$m$, lema titik tetep
nyebut yen ana fungsi komputabel parsial~$f$
kang nyukupi persamaan ing ngisor iki kanggo saben~$y$:
\[
f(y) \simeq
\begin{cases}
k(y) & \text{yen $l(y) = 0$} \\
f(m(y)) & \text{yen ora.}
\end{cases}
\]
Luwih mligi, $f$~diasilake kanthi netepake
\[
g(x,y) \simeq
\begin{cases}
k(y) & \text{yen $l(y) = 0$} \\
\cfind{x}(m(y)) & \text{yen ora}
\end{cases}
\]
banjur nganggo lema titik tetep kanggo
nemokake indeks~$e$ saengga
$\cfind{e}(y) = g(e,y)$.

Minangka tuladha nyata, fungsi
``pambagi bebarengan paling gedhe''
$\fn{gcd}(u,v)$ bisa ditetepake kanthi
\[
\fn{gcd}(u,v) \simeq
\begin{cases}
v & \text{yen $u = 0$} \\
\fn{gcd}(\fn{mod}(v, u), u) & \text{yen ora}
\end{cases}
\]
ing ngendi $\fn{mod}(v, u)$ nuduhake turahan
nalika $v$ dipara dening~$u$. Lema titik
tetep nuduhake yen $\fn{gcd}$ komputabel parsial.
(Sejatine, iki bisa dilebokake ing pola ing
ndhuwur, kanthi $y$ minangka kode pasangan
$\tuple{u, v}$.) Induksi marang~$u$
banjur nuduhake yen sejatine $\fn{gcd}$ total. Cathetan panyunting: ing pasangan input nol,nol, definisi iki menehi asil nol; iki konvensi kanggo pasangan iku.

Mesthi wae, kita bisa nggawe definisi kang
ngrujuk marang awake dhewe kang luwih
rumit tinimbang tuladha-tuladha mau.
Umume basa pamrograman ngidini definisi
fungsi adhedhasar fungsi iku dhewe,
kanthi sawijine cara. Gatekna yen iki
rada kurang nggumunake tinimbang bisa
netepake fungsi adhedhasar sawijining
\emph{indeks} algoritma kang ngitung fungsi
iku, kaya kang diidinake dening teorema
titik tetep kanthi wujud paling umume.

\end{document}
